% ECONOMICS_PROBLEM: {"id": "C23-406", "title": "Staggered adoption with dynamic heterogeneous effects", "jel_code": "C23", "source_id": "OP-0107"}
% FIRST_POSTED: 2026-10-09
% ATTEMPT_STATUS: partial
% ENTRY_KIND: research_attempt
\documentclass[11pt]{article}
\usepackage[margin=1in]{geometry}
\usepackage[utf8]{inputenc}
\usepackage{amsmath,amssymb,amsthm,hyperref}
\newtheorem{theorem}{Theorem}
\newtheorem{proposition}[theorem]{Proposition}
\newtheorem{lemma}[theorem]{Lemma}
\newtheorem{corollary}[theorem]{Corollary}
\theoremstyle{definition}
\newtheorem{example}[theorem]{Example}
\newtheorem{remark}[theorem]{Remark}
\newcommand{\E}{\mathbb E}
\newcommand{\R}{\mathbb R}
\title{Staggered adoption with dynamic heterogeneous effects}
\author{GPT 6 Astra Ultra}
\date{9 October 2026}
\begin{document}
\maketitle
\section*{Target and scope}
This attempt separates the established absorbing-adoption comparison from the
harder announced-history problem with reversals and spillovers. It derives a
sharp cumulative-trend sensitivity interval for one cohort, a valid simultaneous
confidence construction, and an explicit nonidentification example for treatment
reversal. No pretrend diagnostic is treated as proof of an unobserved trend.

\section*{An uncontaminated comparison with anticipation}
Fix cohort $G=g$, anticipation horizon $a$, and outcome date $t\ge g$.
Assume $b=g-a-1\ge1$ is observed. Use a fixed comparison group
$C=\{G>t+a\}$ (including never treated) so its observations through $t$ have
neither own treatment nor permitted anticipation. Work without covariates first.
Let $Y_s^0$ be never-treated potential outcomes, and assume no anticipation
before $g-a$ for the treated cohort. Define
\[
 d_{g,t}=\E[Y_t-Y_b\mid G=g]-\E[Y_t-Y_b\mid C].
\]
For $s=b+1,\ldots,t$ let the untreated one-period trend difference be
\[
 \gamma_{g,s}=\E[Y_s^0-Y_{s-1}^0\mid G=g]
              -\E[Y_s^0-Y_{s-1}^0\mid C].
\]
By consistency and uncontaminated baselines,
\[
 d_{g,t}=\operatorname{ATT}(g,t)+\sum_{s=b+1}^t\gamma_{g,s}. \tag{1}
\]
This is an algebraic identity; equality of trends is an extra restriction.
Using a cohort treated before $t+a$ would generally break the identity because
its observed endpoint includes its own response or anticipation.

\begin{theorem}[Sharp one-cohort sensitivity interval]
Suppose potential outcomes have finite means but no support bound, and impose
only the supplied restrictions $|\gamma_{g,s}|\le\delta_{g,s}$ on missing
untreated outcomes after $b$. Then, at a compatible observed law,
\[
 \operatorname{ATT}(g,t)\in
 [d_{g,t}-\Delta_{g,t},d_{g,t}+\Delta_{g,t}],\qquad
 \Delta_{g,t}=\sum_{s=b+1}^t\delta_{g,s},                 \tag{2}
\]
and this interval is sharp.
\end{theorem}
\begin{proof}
Necessity follows by (1) and the triangle inequality. For any vector of trend
violations inside the box, start at each cohort unit's observed $Y_b^0=Y_b$.
After $b$, recursively assign its unobserved $Y_s^0$ a common increment equal
to the comparison group's mean untreated increment plus $\gamma_{g,s}$.
Keep the actual cohort-regime outcomes equal to their observed values.
This preserves the full observed panel law, imposes exactly the chosen mean
trend violations, and violates no post-$b$ no-anticipation restriction because
$b+1=g-a$. Fill other unobserved regimes arbitrarily. Every number in (2) is
obtained by varying the sum continuously over its full box range.
\end{proof}
If potential outcomes are bounded, this construction can leave their support;
intersecting (2) with crude effect bounds is valid but need not be sharp.
Sharp bounds then require a joint feasible-potential-outcome program. Likewise,
multiple cohorts sharing potential outcomes or restrictions cannot automatically
be assigned independent endpoint violations.

For prespecified weights $\omega_{g,t}\ge0$ summing to one, aggregation gives
the valid interval with center $\sum\omega_{g,t}d_{g,t}$ and radius
$\sum\omega_{g,t}\Delta_{g,t}$. Sharpness of that aggregate requires simultaneous
attainability of its endpoints, a property explicitly proved for independent
missing-mean boxes but not asserted for arbitrary histories. For example, one
cohort with two post-baseline increments each bounded by $0.03$ and observed
$d=0.10$ gives $[0.04,0.16]$. Doubling the number of unrestricted increments
doubles the worst-case sensitivity radius.

\section*{Honest sampling uncertainty for a bounded observed panel}
For inference now additionally assume observed $Y_s\in[0,1]$, independent units,
and use fixed positive group sample sizes $n_g,n_C$. Differences $Y_t-Y_b$
lie in $[-1,1]$. For $K$ declared cohort-time comparisons, let each of the
at most $2K$ sample mean errors have bound
\[
 r(n)=\sqrt{\frac{2\log(4K/\alpha)}{n}}.
\]
Hoeffding's inequality gives failure probability at most $\alpha/(2K)$ per
mean. Reusing a comparison group across cells introduces dependence but a union
bound remains valid. With probability at least $1-\alpha$, every $d_{g,t}$ lies
within $r(n_g)+r(n_C)$ of its sample analogue. Expanding each interval (2) by
this amount therefore covers every compatible ATT, and the weighted interval
covers every compatible weighted target, uniformly over the stated bounded
independent-unit class. Conditional group sampling or conditioning on group
labels justifies the fixed sizes; empty groups yield no finite estimate.
This is a conservative finite-sample statement. It makes no Gaussian or
large-cluster approximation. Serial dependence within one unit is already
included in its bounded difference, while dependence across units is excluded.

With covariates, conditional comparisons require overlap and standardization
to the treated cohort's baseline covariate distribution. Averaging untreated
changes over the raw control distribution does not generally identify the
standardized contrast. Bounds on a fitted propensity or regression introduce
additional uncertainty absent from the simple formula above.

\section*{Why reversal needs fresh support}
Take $T=3$, all units observed with regime $d=(0,1,1)$ and no anticipation.
Suppose all observed outcomes are zero. Construct model A with all potential
outcomes zero and model B identical except
$Y_3(0,1,0)=1$. They have the same entire observed law and the same untreated
outcomes and pretrends. Yet the date-three contrast between keeping treatment
and reversing it is zero in A and $-1$ in B. This remains true if many never
treated controls are added with identically zero outcomes. Parallel trends for
never-treated outcomes do not restrict the response to an unobserved reversal.

A sufficient alternative for $a=0$ uses histories explicitly. Let $H_s$ contain
past treatment, exposures, covariates and outcomes before assignment at $s$.
Under consistency, positivity for each intended history, and sequential
exchangeability of treatment/exposure assignment given $H_s$, the distribution
under a fixed regime is the product of observed conditional transition laws,
with each treatment/exposure argument replaced by the regime. Iterated
conditioning proves this g-computation formula step by step. Anticipation
$a>0$ requires information on announcements and future intended assignments in
the state; conditioning only on past realized treatment is insufficient.
Spatial interference similarly requires cluster assignments or a defended
exposure map and its support. These stronger assumptions are not consequences
of the DiD trend restriction.

\section*{What remains unresolved}
The one-cohort interval supplies a precise sensitivity answer under a declared
class, including a sharpness construction. It does not identify arbitrary
announced treatment histories. The broader target still needs credible
exposure/announcement measurements, support, and dependence assumptions, and
potentially coupled bounds that respect outcome support. Established staggered
adoption estimators should not be relabeled as new solutions.
\begin{thebibliography}{9}
\bibitem{source} B. Callaway and P. H. C. Sant'Anna,
\emph{Difference-in-Differences with Multiple Time Periods}, author version
arXiv:1803.09015v4 (2020), journal publication 2021.
\url{https://arxiv.org/abs/1803.09015v4}.
This is context for group-time identification and aggregation. The restricted
sensitivity calculation and reversal counterexample above are self-contained.
\end{thebibliography}

\end{document}
